\section{MAST}
The response time analysis of the system has been largely performed using the
\textit{MAST} tool \cite{MASTComprehensiveTool}. \textit{MAST} is a
\textit{Modeling and Analysis Suite for Real-Time} applications, developed at
the \textit{University of Canatabria}. It enables the modelling of real-time
systems and provides a set of tools for performing various types of analysis on
a given model. \\ 
In particular, for the purpose of this project, \textit{MAST} can perform
response-time analysis for real-time systems on a single-processor using a
\textit{Fixed-Priority Scheduling} policy. The tool also supports preemptive 
tasks and the \textit{Immediate Ceiling Protocol (ICP)} for managing shared
resources. This latter support stems from the fact that \textit{MAST} was
originally designed to analyze \textit{Ada} systems, which implements the
\textit{ICP} for resource access control. Since \textit{ICP} and \textit{SRP}
exhibit similar observable behaviour, \textit{MAST} can be used to analyze
\textit{SRP}-based system, such as our \textit{Rust} implementation.

\subsection{MAST Model}
To analyze the system using \textit{MAST}, a model of the system must be
provided in the form of an XML or text file. In this sub-section, we describe
the model components available for representing the system. Furtrher details can
be found in the \textit{MAST description} \cite{mastdescription}.

\subsubsection{Processing\_Resource} 
It represents a resource capable of executing abstract activities. In our case,
since we are dealing with a single-processor system, there is only one
processing resource. When we define a processing resource, we must provide the
range of priorities available, a \textit{System\_Timer} and the worst
\textit{ISR Switch} time. The latter parameter denotes the context switch time
of an interrupt service routine, without taking into account the execution time
of the ISR itself.

\subsubsection{System\_Timer} 
It represents a timer used by a processing resource. In our system we have to
deal with \textit{Ticker}, a system that has a periodic ticker. The
informations required by this type of timer are the period and the worst-case
overhead time, i.e., the time taken to process the ticker interrupt, including
the scheduling overhead required for scanning and moving the jobs in the ready
queue. 

\subsubsection{Scheduler}
It represents the scheduling policy used to manage the amount of CPU processing
capacity. As previously mentioned, our system uses a \textit{Fixed-Priority
Scheduling} policy. When defining a scheduler, we need to specify the processing
resource it is hosted on, the available priority range and the worst context
switch time. This last parameter denotes the time taken to set 
the context switch interrupt as pending and to execute its handler, which is
responsible for saving registers of active context and for restoring the ones of
the new active context.

\subsubsection{Shared\_Resource}
It represents a resource that is shared among tasks. As stated before, the
\textit{MAST} model uses the \textit{Immediate Ceiling Protocol}; therefore the
priority specified for a shared resource is its ceiling priority, i.e. the
highest priority of those tasks which access it.

\subsubsection{Scheduling\_Server}
It represents a schedulable entity in a processing resource. Since we use a
\textit{Fixed-Priority Scheduling} policy, the server scheduling parameters are
of type \textit{Fixed\_Priority\_Policy}, which corresponds to regular preemptive
fixed-priority scheduling, and require a priority value to be specified.

\subsubsection{Operation}
It is an abstract representation of a piece of code. There are three types of
operations:
\begin{itemize}
    \item \textit{Simple}: used to model simple operations or methods of
    protected objects. In the latter case, the operation requires a list of
    shared resources to lock before executing it and to unlock after its
    completion. It also requires the specification of best, average and worst-case
    execution times;
    \item \textit{Composite}: an operation composed of an ordered sequence of
    other operations. Its execution times are derived from the execution times
    of its internal operations;
    \item \textit{Enclosing}: similar to a composite one but we need to specify
    the best, average and worst-case execution times as well. This provides a
    way to express overheads associated with inner operations, such as the cost
    of acquiring and releasing a shared resource. 
\end{itemize}

\subsubsection{Transaction} 
It represents a modeling unit for interactions of the system. Each transaction
is composed of three components:
\begin{itemize}
    \item \textit{External\_Events}: a list of events that models interactions
    of the system with external devices through interrupts or with hardware
    timing devices. For each external event it is necessary to specify its type,
    such as periodic or sporadic, its period or minimum inter-arrival time, its
    jitter and its phase;
    \item \textit{Internal\_Events}: a list of events that are generated by
    event handlers. They can be used to express deadlines through
    \textit{Timing\_Requirements}: in our case, we use \textit{Hard\_Global\_Deadline}
    timing requirements, which specify an hard deadline that must be met by all
    instances of the internal event, and a reference to the external event that
    represents its activation;
    \item \textit{Event\_Handlers}: a list of event handlers that enable to
    express relationships between events. Each event handler requires an event
    (or a list of events) as input and produces an event (or a list of events) as
    output. Based on the type of the event handler (such as
    \textit{System\_Timed\_Activity}), it is also possible to execute an activity
    operation on a specified activity server. 
\end{itemize}

\subsection{MAST Tools}
\textit{MAST} supports different types of analysis tools. In this subsection, we
present those used in our system's analysis, while further details
can be found in the \textit{MAST description} \cite{restrictionsMast}.
\begin{itemize}
    \item \textit{Classic RM Analysis}: the classic exact response-time analysis
    for single-processor fixed-priority systems.
    Although the name might suggest support only for \textit{Rate Monotonic}
    scheme, it also covers \textit{Deadline Monotonic} as well, which is the one
    adopted in our system; 
    \item \textit{Offeset-based Approximate with Precedence Relation Analysis}: 
    response-time analysis that improves the holistic analysis by taking into
    account that tasks of the same transaction are not independent, through the
    use of offsets. Those precedence relations, togheter with tasks priorities,
    provide a tighter estimation of the response times comparing to other tools.
\end{itemize}
Each analysis tool, as described in the \textit{MAST restrictions}
\cite{restrictionsMast}, is subject to certain limitations regarding the
expression power of the model, particulary in how transactions can be
constructed. The most relavant are: 
\begin{itemize}
    \item\textit{Restricted\_Multipath\_Transactions}: it imposes the absence of
    branch elements in side transactions;
    \item \textit{Referenced\_Events\_Are\_External\_Only}: it requires only
    external events as referenceable events, therefore only external events can
    be used as activation events for deadlines; 
    \item\textit{Linear\_Transactions\_Only}: it allows just one external event
    per transaction. 
\end{itemize}